% Kilitli üretim tasarımı. Üç panel değil bir panel: bir Latin hiperküpün
% üç izdüşümü de aynı şeyi söyler, o da hiperküpün kutusunu doldurduğudur.
% Gösterilmeye değen şey, iç içe geçmiş alt kümedir.
% `docs/rapor2204` içinden derle:  lualatex sekil/tasarim.tex
\documentclass[border=3pt,tikz]{standalone}
\input{sekil/ortak.tex}

\colorlet{ink}{black!82}
\colorlet{incecol}{red!72!black}

\pgfplotstableread{\VERI/tasarim_kaba.dat}\kaba
\pgfplotstableread{\VERI/tasarim_ince.dat}\ince

\begin{document}
\begin{tikzpicture}[font=\footnotesize]
% Eksen, ön baskıdakinden (62x46 mm) büyük: rapor tek kolon olduğu için
% şekil sayfada ~2 kat büyütülüyordu ve yazılar gövde metnine göre
% orantısız kalıyordu. Doğal boyut büyüyünce büyütme oranı ~1,2'ye düşer.
\begin{axis}[
  width=92mm, height=66mm,
  ymode=log,
  xmin=0.985, xmax=1.315, ymin=0.6, ymax=1.7e5,
  xlabel={blok gözenekliliği $\alpha_b$},
  ylabel={matris dayanımı $Y_0$\:[Pa]},
  xtick={1.0,1.1,1.2,1.3},
  xticklabels={$1{,}0$,$1{,}1$,$1{,}2$,$1{,}3$},
  ytick={1,100,10000}, yticklabels={$1$,$10^2$,$10^4$},
  axis line style={draw=ink!65}, tick style={draw=ink!65},
  ticklabel style={scale=0.85, text=ink!80},
  label style={scale=0.92, text=ink},
  xlabel style={below=-2pt}, ylabel style={above=-4pt},
  axis on top=false,
]
  \addplot[only marks, mark=o, mark size=1.2pt,
           draw=ink!35, line width=0.35pt]
    table[x=alphab, y=Y0] {\kaba};
  \addplot[only marks, mark=*, mark size=1.7pt,
           draw=incecol, fill=incecol]
    table[x=alphab, y=Y0] {\ince};
\end{axis}

\begin{scope}[shift={($(current axis.north west)+(1mm,5.5mm)$)}]
  \draw[ink!35, line width=0.35pt] (0,0) circle (1.2pt);
  \node[anchor=west, text=ink, font=\scriptsize, inner sep=3pt] at (1.6mm,0)
       {84 yalnız kaba};
  \fill[incecol] (40mm,0) circle (1.7pt);
  \node[anchor=west, text=ink, font=\scriptsize, inner sep=3pt] at (41.6mm,0)
       {12 iç içe ince};
\end{scope}

\end{tikzpicture}
\end{document}
